% =====================================================================
%  MATH300 Project Description -- Discrete Maps
% =====================================================================
\documentclass[11pt,a4paper]{article}

\usepackage[T1]{fontenc}
\usepackage[utf8]{inputenc}
\usepackage{lmodern}
\usepackage[margin=2.5cm]{geometry}
\usepackage{enumitem}
\usepackage{booktabs}
\usepackage{tabularx}
\usepackage{array}
\usepackage{amsmath,amssymb}
\usepackage{xcolor}
\usepackage[colorlinks=true,urlcolor=blue,linkcolor=black,citecolor=black]{hyperref}

\setlist{itemsep=2pt, topsep=4pt}
\setlength{\parindent}{0pt}
\setlength{\parskip}{6pt}

% Placeholder macro: blue text in square brackets.
% To hide the colour in the final version, change {blue} to {black}.
\newcommand{\ph}[1]{\textcolor{blue}{[#1]}}

% Marker for optional bonus tasks.
\newcommand{\bonus}{\textbf{(Bonus)}}

\newcolumntype{L}{>{\raggedright\arraybackslash}X}

\begin{document}

% ---------------------------------------------------------------------
%  Title (no logo, no instructor name)
% ---------------------------------------------------------------------
\begin{center}
  {\large\bfseries MATH300: Mathematical Computing \textperiodcentered{} Fall 2026}\\[6pt]
  {\Large\bfseries Project: Chaos in Discrete Maps}
\end{center}
\vspace{6pt}

Groups must be formed, and project titles submitted on Canvas, by
\textbf{October 26, 2026}. The project
has two graded deliverables and an oral exam: a 3--5 page report for
Project~2 (due \textbf{November 20, 2026}), a poster for the final project
(due \textbf{December 11, 2026}), and an oral exam held between November 30
and December 17, 2026, 3:30\,pm. Read the general rules below before you
start; they apply throughout the project.

% ---------------------------------------------------------------------
\section*{General remarks}
% ---------------------------------------------------------------------
\begin{itemize}
  \item Projects are set for a group of \textbf{two or three} students.
  \item Form your group by \textbf{October 26, 2026}, and submit your
        project title as a group assignment on \textbf{Canvas} by the same
        date.
  \item Every group member must carry an \textbf{equal portion} of the coding
        and of the report and poster preparation.
  \item The scope can be adjusted to the number of participants, their
        interests and prior knowledge.
  \item Only \textbf{Python} may be used for programming.
  \item Implementation, documentation, validation and demonstration examples
        are all part of the project.
  \item Within these rules you are free to choose the implementation, code
        structure and exact content. You can discuss this with me at any time.
\end{itemize}

% ---------------------------------------------------------------------
\section*{Version control and GitHub}
% ---------------------------------------------------------------------
Each group keeps \textbf{one GitHub repository for the whole project}
(Project~2 and the final project), and every member commits from their own
GitHub account.
\begin{itemize}
  \item Use git from the very beginning. Commits should be small, consistent
        and clearly described. Avoid large commits.
  \item Each group member uploads their own code using \textbf{their own
        GitHub account}. Commits from a shared or borrowed account will not
        be credited to the other members.
  \item The repository name is up to your group.
  \item If the repository is private, add my GitHub account
        \textbf{\texttt{say-yas}} as a collaborator.
  \item The repository must contain a README describing the project, how to
        install and run the code, and the AI-use statement (see below).
\end{itemize}

% ---------------------------------------------------------------------
\section*{Use of AI tools}
% ---------------------------------------------------------------------
\textbf{AI-generated code is not accepted in any form.} Code produced by a
chat assistant or code generator (e.g.\ ChatGPT, Claude, Gemini) and copied,
pasted or adapted into your project will be treated as work that is not your
own.
\begin{itemize}
  \item \textbf{Permitted:} inline coding assistance such as GitHub Copilot
        autocompletion while you write the code yourself, \textbf{provided its
        use is acknowledged}.
  \item \textbf{Acknowledgement:} include an AI-use statement in the README
        and in the report that names the tool(s) used and the parts of the
        code (files or functions) where they assisted. If no AI tools were
        used, state this explicitly.
  \item \textbf{Reports and posters must not be generated by AI} in any
        form: their text, figures and layout must be your own work. If AI
        tools were used for any other non-code purpose (e.g.\ checking
        grammar or searching the literature), describe this in the same
        statement.
  \item You are responsible for every line in your repository. In the oral
        exam you must be able to explain and modify any part of the code you
        submitted.
\end{itemize}

% ---------------------------------------------------------------------
\section*{Deliverables and deadlines}
% ---------------------------------------------------------------------
Submit each deliverable following the details provided on the
\textbf{Canvas} page. The written report and poster are handed in on
\textbf{Crowdmark}, and the link to your GitHub repository is submitted
separately on \textbf{Canvas}.

\begin{center}
\small
\renewcommand{\arraystretch}{1.35}
\begin{tabularx}{\textwidth}{|>{\raggedright\arraybackslash}p{2.5cm}|>{\raggedright\arraybackslash}p{3.0cm}|>{\raggedright\arraybackslash}p{3.0cm}|L|}
\hline
\textbf{Deliverable} & \textbf{Due} & \textbf{Submission} & \textbf{Notes} \\
\hline\hline
Group and project title & Monday, October 26, 2026 &
  Group assignment on Canvas &
  Group of two or three members; submit the project title. \\
\hline
Project~2 report & Friday, November 20, 2026 &
  Report (PDF) on Crowdmark; repository link on Canvas &
  3--5 pages; the page limit \textbf{includes} references, figures and
  appendices. \\
\hline
Final project poster & Friday, December 11, 2026 &
  Poster (PDF) on Crowdmark; repository link on Canvas &
  36 inches high by 48 inches wide (landscape). Target audience:
  undergraduate classmates with a basic background in math, physics or
  engineering. \\
\hline
Oral exam & Between November 30 and December 17, 2026, 3:30\,pm &
  Individual meeting; details discussed in class &
  30 minutes per student. Your poster must already be submitted on
  Crowdmark before your exam. The last possible slot ends December 17 at
  3:30\,pm. \\
\hline
\end{tabularx}
\end{center}

\subsection*{Oral exam}
The oral exam is individual and lasts 30 minutes per student. It takes place
between November 30 and December 17, 2026, 3:30\,pm. The oral exams and how
to arrange them will be discussed in class before the Thanksgiving break.

Before your oral exam, your poster must also be submitted on
\textbf{Crowdmark}; that is, you must have finished both Project~2 and the
final project. In the meeting you will:
\begin{itemize}
  \item Present your results and algorithms.
  \item Possibly live-code and run parts of your algorithms to generate
        outputs on the spot.
  \item Answer questions about the theory, your implementation and your
        results. Each member should be able to explain the whole project, not
        only their own parts.
\end{itemize}
Bring a laptop with the repository cloned and a working environment, so that
the code runs without additional setup.

% ---------------------------------------------------------------------
\section*{Grading}
% ---------------------------------------------------------------------
\begin{center}
\small
\renewcommand{\arraystretch}{1.35}
\begin{tabularx}{0.8\textwidth}{|L|>{\raggedleft\arraybackslash}p{3.2cm}|}
\hline
\textbf{Component} & \textbf{Points} \\
\hline\hline
Group and project title (on time) & 10 \\
\hline
Project~2 report & 100 \\
\hline
Final project poster & 100 \\
\hline
Oral exam (individual) & 50 \\
\hline
Bonus: extra steps beyond the required tasks & up to 10 \\
\hline
\end{tabularx}
\end{center}

The Project~2 report and the poster are each graded out of 100. What matters
most is:
\begin{itemize}
  \item \textbf{Coherence:} a clear story from the question, through the
        methods, to the results and conclusions.
  \item \textbf{Correctness:} of the mathematics, the code and the results.
  \item \textbf{Depth of analysis and discussion:} explain what your results
        mean, why they look the way they do, and what their limitations are,
        rather than only showing them.
  \item \textbf{Professional presentation:} figures with labeled axes and
        readable font sizes, captions and legends, a clean and readable
        layout, and proper references.
\end{itemize}

A complete, correct and well-discussed
solution of the required tasks can earn full marks. Tasks marked
\bonus{} and the suggestions under ``Going further'' are optional; exploring
such extra steps can earn \textbf{up to 10 bonus points} in total.

% ---------------------------------------------------------------------
\section*{Requirements for the report and the poster}
% ---------------------------------------------------------------------
Both the report and the poster must contain a clear description of the
project, references, properly labeled figures and an author-contributions
statement.

\subsection*{Content}
\begin{itemize}
  \item A clear and concise description of the problem, the methods and the
        results.
  \item The AI-use statement described above.
  \item The report should not contain any significant amount of code; refer
        to the repository instead.
\end{itemize}

\subsection*{Poster}
\begin{itemize}
  \item The poster must be \textbf{self-contained}: a reader should be able to
        follow it without the report, the code or a spoken explanation.
  \item You may be selective in what you show and how you narrate your story,
        so that an undergraduate audience with a basic background in math,
        physics or engineering can follow it.
  \item You may choose your own title for the poster; be creative.
\end{itemize}

\subsection*{References}
\begin{itemize}
  \item Include a list of references at the end of the report and the poster.
  \item Cite sources clearly where they are used in the text (papers, books,
        software, code repositories and data sets).
  \item Use one consistent citation style throughout.
\end{itemize}

\subsection*{Figures}
\begin{itemize}
  \item Every plot must have labeled $x$ and $y$ axes (and $z$ axis for 3D
        plots or color bars), including units where applicable.
  \item Axis labels, tick labels and legends must have a font size similar to
        the surrounding text, so they remain readable without zooming.
  \item Each figure needs a legend when it shows more than one data series,
        and a caption that states what is shown.
  \item Use log scales where appropriate (e.g.\ for convergence plots) and say
        so in the caption.
\end{itemize}

\subsection*{Author contributions}
Each group must include an author-contributions statement in the report and
the README that states:
\begin{itemize}
  \item Who did what: which parts of the code, experiments, analysis, writing
        and poster each member worked on.
  \item To what extent: the approximate share of each member in each part
        (e.g.\ 50\,\% / 50\,\%).
  \item Who took the lead in which parts.
\end{itemize}
The coding and the report and poster preparation must be shared equally
between the members. The commit history should be consistent with the stated
contributions.

% ---------------------------------------------------------------------
\section*{Codebase requirements}
% ---------------------------------------------------------------------
\begin{itemize}
  \item Well structured and documented, both in the code (docstrings,
        comments) and in the repository (README).
  \item Validated: show that your implementation produces correct results,
        e.g.\ on examples with known solutions.\end{itemize}

% ---------------------------------------------------------------------
\section*{Project description}
% ---------------------------------------------------------------------

\subsection*{Goal}
Very simple rules can produce very complicated behavior. In this project you
will study \emph{discrete maps}: rules that take the current state of a
system and give its state at the next time step. You will choose
\textbf{at least 3 maps} whose time domain is discrete from the list in
\cite{list}, implement them in Python, and explore their behavior
numerically: their trajectories, how they depend on the parameters, and when
they become chaotic.

\subsection*{Background}
A discrete map is a function $f$ that is applied again and again,
\begin{equation*}
  x_{k+1} = f(x_k), \qquad k = 0, 1, 2, \dots,
\end{equation*}
starting from an initial state $x_0$. The state can be a single number or a
point in the plane. The sequence $x_0, x_1, x_2, \dots$ is called a
\emph{trajectory} (or orbit). A classical example is the logistic map
\cite{may},
\begin{equation*}
  x_{k+1} = r\,x_k(1 - x_k), \qquad 0 \le x_k \le 1,
\end{equation*}
a simple model for population growth. For small values of the parameter $r$
the trajectory settles to a fixed value; as $r$ increases it starts to
oscillate between 2, then 4, then 8 values (period doubling), and for larger
$r$ it becomes \emph{chaotic}: it never repeats, and two trajectories that
start very close together soon become completely different. A two-dimensional
example is the H\'enon map \cite{henon},
\begin{equation*}
  x_{k+1} = 1 - a\,x_k^2 + y_k, \qquad y_{k+1} = b\,x_k,
\end{equation*}
whose trajectories are attracted to a fractal-shaped set, a \emph{strange
attractor}.

Two simple pictures help to see the structure hidden in a trajectory.

The \emph{phase space} of a map is the set of all its possible states: the
interval $[0, 1]$ for the logistic map, or the plane of all points $(x, y)$
for the H\'enon map. A \emph{phase portrait} is a picture of trajectories
drawn in phase space. Time is not shown on any axis; instead, each state
$(x_k, y_k)$ of a trajectory is drawn as one point, so the trajectory appears
as a set of dots. For a two-dimensional map you plot $y_k$ against $x_k$ for
all $k$, often for several initial states and usually after discarding the
first iterations (the transient). The picture then shows where trajectories
end up in the long run: a single point (a stable fixed point), a few points
visited in turn (a periodic orbit), a closed curve, or a scattered,
fractal-looking set (a strange attractor, as for the H\'enon map). For maps
without an attractor, such as the standard map, it shows which regions are
regular and which are chaotic.

A \emph{time-delay embedding} does the same with a single sequence: it
plots $x_{k+\tau}$ against $x_k$ for a delay $\tau$ (for example $\tau = 1$),
or the three-dimensional points $(x_k, x_{k+\tau}, x_{k+2\tau})$. A
trajectory that looks random in the time series often turns out to lie on a
simple curve or a fractal set in such a plot. A good introduction to all of
these ideas is \cite{strogatz}.

\subsection*{Tasks for Project 2 (report, due November 20, 2026)}
\begin{enumerate}
  \item Choose \textbf{at least 3 maps} with a discrete time domain from
        \cite{list}. Include at least one one-dimensional map (such as the
        logistic map or the tent map) and at least one two-dimensional map
        (such as the H\'enon map or the standard map). Describe each map and
        the meaning of its parameters, and cite your sources.
  \item Implement a general function that iterates a given map from an
        initial state for a given number of steps, and implement your chosen
        maps with it.
  \item For each map, plot trajectories for a few parameter values: the time
        series $x_k$ against $k$, and a phase portrait (for the
        two-dimensional maps) or a time-delay embedding (for the
        one-dimensional maps).
  \item \bonus\ Make three-dimensional time-delay embeddings
        $(x_k, x_{k+\tau}, x_{k+2\tau})$ and compare different delays
        $\tau$.
  \item \bonus\ Study more than 3 maps.
\end{enumerate}

\subsection*{Tasks for the final project (poster, due December 11, 2026)}
\begin{enumerate}
  \item For each map, choose one parameter and show how the behavior changes
        as you vary it: make a sequence of phase portraits (two-dimensional
        maps) or time-delay embeddings (one-dimensional maps) for several
        parameter values, from regular to chaotic behavior. Describe what
        changes.
  \item Discuss what you observe: where is each map regular and where is it
        chaotic, and what do the maps have in common? What are the
        limitations of studying chaos numerically (finite precision,
        transients, finite number of iterations)?
  \item \bonus\ For a few parameter values in the chaotic range, plot a
        histogram of the values visited by a long trajectory. Which values
        does the trajectory visit most often?
  \item \bonus\ Estimate the parameter value at which the behavior becomes
        chaotic, and explain how you found it.
\end{enumerate}

\subsection*{Going further}
These optional steps count toward the bonus points. For example, you can draw
a \emph{bifurcation diagram}: for each value of a parameter (horizontal
axis), iterate the map, discard the transient, and plot the values the
trajectory visits afterwards. It shows period doubling as branches that
split in two, and chaos as a cloud of dots. You can also study more maps
from the list.

% ---------------------------------------------------------------------
\begin{thebibliography}{9}
  \bibitem{list} Wikipedia. \emph{List of chaotic maps}.
    \url{https://en.wikipedia.org/wiki/List_of_chaotic_maps}
  \bibitem{may} R.~M.~May (1976). Simple mathematical models with very
    complicated dynamics. \emph{Nature}, 261, 459--467.
    \url{https://doi.org/10.1038/261459a0}
  \bibitem{henon} M.~H\'enon (1976). A two-dimensional mapping with a strange
    attractor. \emph{Communications in Mathematical Physics}, 50, 69--77.
    \url{https://doi.org/10.1007/BF01608556}
  \bibitem{strogatz} S.~H.~Strogatz (2015). \emph{Nonlinear Dynamics and
    Chaos}, 2nd edition. Westview Press.
\end{thebibliography}

\end{document}
